% from ega0-1-fr.tex
\oldpage[0\textsubscript{I}]{60}

ouvert $U$ de $y$ et un faisceau cohérent d'idéaux $\mathcal J$ de
$\mathcal O_Y|U$ tel que $\mathcal J_y=\mathfrak J$, d'où la conclusion.
\end{env}

\begin{env}[6.7.4]
\label{0.6.7.4-fr}
Les résultats de (6.1) sur les modules plats se transcrivent aussitôt en
propositions sur les faisceaux $f$-plats en un point~:

Si $0\to\mathcal F'\to\mathcal F\to\mathcal F''\to0$ est une suite exacte
de $\mathcal O_X$-Modules et si $\mathcal F''$ est $f$-plat au point
$x\in X$, alors, pour tout voisinage ouvert $U$ de $y=f(x)$ et tout
$(\mathcal O_Y|U)$-Module $\mathcal G$, la suite
\[
  0\longrightarrow
  (f^*(\mathcal G)\otimes_{\mathcal O_X}\mathcal F')_x
  \longrightarrow
  (f^*(\mathcal G)\otimes_{\mathcal O_X}\mathcal F)_x
  \longrightarrow
  (f^*(\mathcal G)\otimes_{\mathcal O_X}\mathcal F'')_x
  \longrightarrow0
\]
est exacte. Pour que $\mathcal F$ soit $f$-plat en $x$, il faut et il suffit
alors que $\mathcal F'$ le soit. On a des énoncés analogues pour les notions
correspondantes de $\mathcal O_X$-Module $f$-plat au-dessus de $y\in Y$, ou
de $\mathcal O_X$-Module $f$-plat.
\end{env}

\begin{env}[6.7.5]
\label{0.6.7.5-fr}
Soient $f:X\to Y$, $g:Y\to Z$ deux morphismes d'espaces annelés~; soient
$x\in X$, $y=f(x)$, $\mathcal F$ un $\mathcal O_X$-Module. Si $\mathcal F$
est $f$-plat au point $x$ et si le morphisme $g$ est plat au point $y$, alors
$\mathcal F$ est $(g\circ f)$-plat en $x$ (6.2.1). En particulier, si $f$ et
$g$ sont des morphismes plats, $g\circ f$ est plat.
\end{env}

\begin{env}[6.7.6]
\label{0.6.7.6-fr}
Soient $X$, $Y$ deux espaces annelés, $f:X\to Y$ un morphisme plat. Alors
l'homomorphisme canonique de bifoncteurs (4.4.6)
\[
  f^*(\mathcal{H}om_{\mathcal O_Y}(\mathcal F,\mathcal G))
  \longrightarrow
  \mathcal{H}om_{\mathcal O_X}(f^*(\mathcal F),f^*(\mathcal G))
  \tag{6.7.6.1}
  \label{0.6.7.6.1-fr}
\]
est un isomorphisme lorsque $\mathcal F$ admet une présentation finie
(5.2.5).

En effet, la question étant locale, on peut supposer qu'il existe une suite
exacte $\mathcal O_Y^m\to\mathcal O_Y^n\to\mathcal F\to0$. Or, les deux
membres de (6.7.6.1) sont des foncteurs exacts à gauche en $\mathcal F$ en
vertu de l'hypothèse sur $f$~; on est alors ramené à démontrer la proposition
lorsque $\mathcal F=\mathcal O_Y$, cas où elle est triviale.
\end{env}

\begin{env}[6.7.8]
\label{0.6.7.8-fr}
On dit qu'un morphisme $f:X\to Y$ d'espaces annelés est fidèlement plat si
$f$ est surjectif et si, pour tout $x\in X$, $\mathcal O_x$ est un
$\mathcal O_{f(x)}$-module fidèlement plat. Lorsque $X$ et $Y$ sont des
espaces annelés en anneaux locaux (5.5.1), il revient au même de dire que le
morphisme $f$ est surjectif et plat (6.6.2). Lorsque $f$ est fidèlement plat,
$f^*$ est un foncteur exact et fidèle dans la catégorie des
$\mathcal O_Y$-Modules (6.6.1, $a)$) et pour qu'un $\mathcal O_Y$-Module
$\mathcal G$ soit $Y$-plat, il faut et il suffit que $f^*(\mathcal G)$ le
soit (6.6.3).
\end{env}

\section{Anneaux adiques}
\label{section:0.7-fr}

\subsection{Anneaux admissibles}
\label{subsection:0.7.1-fr}

\begin{env}[7.1.1]
\label{0.7.1.1-fr}
Rappelons que dans un anneau topologique $A$ (non nécessairement séparé), on
dit qu'un élément $x$ est topologiquement nilpotent si $0$ est une limite de
la suite $(x^n)_{n\geq0}$. On dit qu'un anneau topologique $A$ est
linéairement topologisé s'il existe un système fondamental de voisinages de
$0$ dans $A$ formé d'idéaux (nécessairement ouverts).
\end{env}

\begin{definition}[7.1.2]
\label{0.7.1.2-fr}
Dans un anneau linéairement topologisé $A$, on dit qu'un idéal $\mathfrak J$
est un idéal de définition si $\mathfrak J$ est ouvert et si, pour tout
voisinage $V$ de $0$, il existe un entier $n>0$ tel

